\scanpage{77}
$\mathcal M$ \deleted{$\unsure{\xi}$} cat. $k$-abélienne avec
\deleted{suffis} $\varinjlim$ et suff. en projectifs
\inserted{et syst. gén.}

$1_M$ tel que
\[
\mathrm{Ab}_k\longrightarrow\mathcal M\qquad
\kappa^*: \xi\longmapsto1_M\otimes_k\xi
\]
commute aux \deleted{produits} \inserted{$\varprojlim$}
i.e. admet un \uline{adjoint à gauche} $\kappa_!$
(alors qu'il admet un adj. à droite
$\kappa_*:a\longmapsto\operatorname{Hom}(1_M,a)$).
\deleted{Conditions}

Alors $\kappa_!$ transforme projectifs en projectifs.

\deleted{Supposons $\mathcal M=M^{\unsure{\flat|b}}$, alors les}

$\kappa^*$ foncteurs \inserted{$:\mathrm{Ab}_k\longrightarrow\mathcal M$, $k$-add}
qui ont un adjoint à gauche et à droite

i.e. qui \deleted{\unsure{ont} commutent aux $\varprojlim$}
commutent aux $\varprojlim$ et $\varinjlim$

Objets $1_M$ tels que

foncteurs $\kappa_!:\mathcal M\mathbin{\text{---}}\mathrm{Ab}_k$
qui commutent aux $\varinjlim$, et

dont l'adjoint \inserted{à dr.} commute aux $\varinjlim$ aussi.

Ex. $\mathcal M=M^{\unsure{\flat|b}}$:
foncteur $M\xrightarrow{\kappa_M}\mathrm{Ab}_k$ qui se factorise par
$\unsure{\mathrm{fab}_k}$

\uline{NB} $1_M:M^o\longrightarrow\mathrm{Ab}_k$
\deleted{\unsure{bien défini $\kappa_!$}}
\inserted{se factorise par $\mathrm{ab}_k$}

soit $\kappa'_M:M^o\longrightarrow\mathrm{ab}_k$,
ce dernier ``n'est autre que'' $\kappa_M^o$

Soit $x$ dans $M$, a-t-on
\[
x\longrightarrow\underbrace{\kappa^*\kappa_!x}_{1_M\otimes\kappa(x)}
\quad\text{épi}
\]
i.e. $\forall y$ dans $M^{\unsure{\flat|b}}$, a-t-on
\[
\begin{array}{c}
\operatorname{Hom}(1_M\otimes\kappa(x),\unsure{y|z})
\longrightarrow\operatorname{Hom}(x,z)\quad\text{mono}\\
\rotatebox{90}{$=$}\\
\operatorname{Hom}_k(\kappa(x),\operatorname{Hom}(1_M,z))\\
\rotatebox{90}{$=$}\\
\operatorname{Hom}(x,\underbrace{1_M\otimes\operatorname{Hom}(1_M,z)}_{f^*f_*z})
\end{array}
\]
\margin{$f^*f_*z\longrightarrow z\quad\text{mono ?}$}

\uline{Ex} $\mathcal M=\mathrm{Ab}_k\qquad1_M=X$ un $k$-module
la condition de commutation \inserted{de $X\otimes_k$} aux $\varprojlim$
signifie
$X$ plat de prés. finie i.e. $X$ proj. de t. f.
\[
M=\mathrm{ab}_k
\]
\[
\kappa^*(\xi)=X\otimes\xi,\qquad\kappa_!(x)=a\otimes\check X
\]
\[
x\longrightarrow\kappa^*(\kappa_!x)=a\otimes\check X\otimes X
\qquad\deleted{\unsure{donc}}\quad\text{épi}\Longleftrightarrow X\text{ de rg }\leq1
\]
\[
\kappa_!\kappa^*(\xi)=\check X\otimes X\otimes\xi\longrightarrow\xi
\qquad\text{mono}\Longleftrightarrow X\text{ de rg }\leq1
\]
\[
\kappa^*\kappa_*x=X\otimes\check X\otimes x\longrightarrow x
\qquad\text{mono}\Longleftrightarrow X\text{ de rg }\leq1
\]
\[
\xi\mathbin{\text{---}}\kappa_*\kappa^*\xi=\check X\otimes X\otimes\xi
\]
\mspage{76}
